% sections/appendix.tex
\section{The full upload ledger}
\label{app:ledger}

\Cref{tab:uploads} lists every scored upload (\val{n-scored} in total; \val{n-rehearsed} with a rehearsal). Dates
are the upload day in US Central time, and a tilde marks a day reconstructed from the campaign ledger. Two
scores from runs that the organisers rescored, one each for \upl{K44} and \upl{F6}, are excluded
(\cref{sec:traps}). \Cref{fig:steps} shows the rehearsal step counts.

\begin{table}[!htbp]
  \centering
  \caption{\textbf{Every scored upload} (\texttt{results/official-scores.csv}). Rehearsal and official in bpb
    (official to four decimals; \upl{K82s4} and \upl{K82s7} were published to seven, \val{best-official-long} and
    \val{salt7-official-long}). Offset $=$ official $-$ rehearsal, rounded half-up from the exact difference, so it
    can differ in the last digit from the offset column of the CSV (\upl{K15}). $\bullet$: new best at the
    time. \textsuperscript{s}\,salt chosen by a lottery; \textsuperscript{i}\,salt inherited from a lottery winner;
    \textsuperscript{n}\,rehearsal normalised to chip-C speed. \upl{K65a}--\upl{K65d} were official tests and were
    not rehearsed. The change descriptions spell out the ledger's internal shorthand.}
  \label{tab:uploads}
  % GENERATED by paper/analysis -- every scored upload, results/official-scores.csv (oracle.py)
{\scriptsize\setlength{\tabcolsep}{3pt}
\begin{tabular}{@{}llcrrrrcp{5.0cm}@{}}
\toprule
Upload & Date & Chip & Rehearsal & Steps & Official & Offset & Best & What changed \\
\midrule
F2 & 09-15 & A & 1.14032 & 1{,}257 & 1.1463 & \ensuremath{+}0.0060 & $\bullet$ & first upload: 131k batch, microbatch 8, cooldown 0.20 \\
F6 & 09-15 & A & 1.13444 & 1{,}316 & 1.1397 & \ensuremath{+}0.0053 & $\bullet$ & cooldown 0.30 \\
K15 & 09-17 & A & 1.12825 & 1{,}808 & 1.1360 & \ensuremath{+}0.0078 & $\bullet$ & scalar parameters as 0-dim tensors, compile, loss computed in a static graph, asynchronous host synchronisation \\
K24 & 09-18 & A & 1.08260 & 912 & 1.0885 & \ensuremath{+}0.0059 & $\bullet$ & depth 12, 262k batch, Muon x2.0, AdamW x1.4142, WD 0.02, cooldown 0.3 \\
K29 & 09-19 & A & 1.04959 & 935 & 1.0554 & \ensuremath{+}0.0058 & $\bullet$ & stream rows (rows laid out as the evaluation reads them) \\
K33 & 09-19 & A & 1.03622 & 917 & 1.0427 & \ensuremath{+}0.0065 & $\bullet$ & value embeddings + stream rows \\
K36 & $\sim$09-21 & A & 1.02058 &  & 1.0268 & \ensuremath{+}0.0062 & $\bullet$ & GQA with one KV head \\
K38 & 09-21 & A & 1.00015 & 1{,}526 & 1.0055 & \ensuremath{+}0.0054 & $\bullet$ & head dim 256 + scalar stack \\
K39 & 09-21 & B & 0.99566 & 1{,}599 & 1.0027 & \ensuremath{+}0.0070 & $\bullet$ & fused CE, time target 1,780 s \\
K40 & 09-21 & A & 0.98954 & 1{,}393 & 0.9953 & \ensuremath{+}0.0058 & $\bullet$ & depth 10 x width 1024, cooldown 0.45, Muon shard; first sub-1.0 \\
K41 & $\sim$09-22 & A & 0.98827 & 1{,}507 & 0.9937 & \ensuremath{+}0.0054 & $\bullet$ & depth 9 x width 1024 \\
K43 & 09-22 & B & 0.98321 & 1{,}610 & 0.9897 & \ensuremath{+}0.0065 & $\bullet$ & seed 67, two EMA/checkpoint fallback levels, EMA every 8 \\
K44 & $\sim$09-24 & B & 0.98235 & 1{,}615 & 0.9888 & \ensuremath{+}0.0065 & $\bullet$ & cooldown 0.60 \\
K45 & 09-25 & B & 0.97941 & 1{,}692 & 0.9861 & \ensuremath{+}0.0067 & $\bullet$ & scalar constants, async optimizer, Muon views \\
K46 & 09-25 & B & 0.97725 & 1{,}753 & 0.9845 & \ensuremath{+}0.0073 & $\bullet$ & fused optimizer update, level 2 \\
K50 & $\sim$09-26 & B & 0.96914 & 1{,}943 & 0.9760 & \ensuremath{+}0.0069 & $\bullet$ & Muon schedule, multi-token prediction, batch warm-up 131k $\rightarrow$ 262k \\
K51 & 09-26 & C & 0.96789 & 2{,}036 & 0.9742 & \ensuremath{+}0.0063 & $\bullet$ & key offset, accumulation schedule 2:0.2,4 \\
K53 & 09-26 & C & 0.96515 & 2{,}142 & 0.9712 & \ensuremath{+}0.0061 & $\bullet$ & ZeRO Muon, EMA prewarm, EMA every 32 \\
K54b\_min & $\sim$09-27 & C & 0.96326 & 2{,}353 & 0.9705 & \ensuremath{+}0.0072 & $\bullet$ & ZeRO-2 + AdamW ZeRO + 3-phase batch ramp \\
K56 & 09-28 & C & 0.96245 & 2{,}360 & 0.9693 & \ensuremath{+}0.0069 & $\bullet$ & key offset radius 0.3, accumulation inside the graph \\
K57 & 09-28 & C & 0.96193 & 2{,}359 & 0.9686 & \ensuremath{+}0.0067 & $\bullet$ & seed 73 \\
K59 & $\sim$09-29 & C & 0.96092 & 2{,}357 & 0.9671 & \ensuremath{+}0.0062 & $\bullet$ & LeakyReLU(0.35)\textsuperscript{2} \\
K60 & 09-29 & C & 0.95910 & 2{,}361 & 0.9655 & \ensuremath{+}0.0064 & $\bullet$ & attention-source reuse 5:6,7,8 \\
K63 & 09-30 & C & 0.95821 & 2{,}385 & 0.9647 & \ensuremath{+}0.0065 & $\bullet$ & EMA off, time target 1,795 s, cooldown 0.7, small-batch LR 0.35/0.6 \\
K65a & 09-30 &  &  &  & 0.9681 &  &  & official test: attention-input shift (XSHIFT) in blocks 6-8 (re-rolled the warm-up path) \\
K65b & 09-30 &  &  &  & 0.9652 &  &  & official test: late k=8 batch in the cooldown \\
K65c & 09-30 &  &  &  & 0.9647 &  &  & official test: EMA-Nesterov lookahead 0.3 (tie) \\
K65d & 09-30 &  &  &  & 0.9653 &  &  & official test: split AdamW cooldown 0.4 \\
K70a\textsuperscript{n} & 09-30 & E & 0.95595 & 2{,}318 & 0.9625 & \ensuremath{+}0.0066 & $\bullet$ & row pool 256; rehearsal normalised to chip-C speed (raw chip E 0.957259, chip E runs 3.5\% slow) \\
K73s4\textsuperscript{s} & 09-30 & H & 0.95521 & 2{,}390 & 0.9620 & \ensuremath{+}0.0068 & $\bullet$ & pool + fused optimizer update, level 3 + execution-queue depth 48 + MLP output-projection LR x0.8, shuffle salt 4 \\
K73s6\textsuperscript{s} & 09-30 & G & 0.95514 & 2{,}398 & 0.9625 & \ensuremath{+}0.0074 &  & shuffle salt 6 (its local edge did not transfer) \\
K77a\textsuperscript{i} & 09-30 & G & 0.95473 & 2{,}375 & 0.9617 & \ensuremath{+}0.0070 & $\bullet$ & K73s4 + EMA blend 0.6 (best official before the final uploads) \\
K82s4\textsuperscript{i} & 09-30 & H & 0.95447 & 2{,}388 & 0.9614 & \ensuremath{+}0.0069 & $\bullet$ & K77a + EMA prewarm; projected 0.9614-0.9615 (range 0.9611-0.9619); best official of the campaign \\
K82s7\textsuperscript{s} & 10-01 & G & 0.95438 & 2{,}405 & 0.9614 & \ensuremath{+}0.0070 &  & K82s4 with shuffle salt 7 (best rehearsal of the campaign); projected about 0.9614 \\
\bottomrule
\end{tabular}
}

\end{table}

\begin{figure}[!htbp]
  \centering
  \includegraphics{step_counts.pdf}
  \caption{\textbf{Optimizer steps in the 30-minute budget}, from the rehearsal of each upload (cold, but warm for
    \upl{K15} and \upl{K38}); the letter is the chip. Shading marks the batch regime: a 131k-token batch before
    \upl{K24}, a constant 262k batch from \upl{K24}, and from \upl{K50} a batch warm-up through 65k and 131k phases,
    which adds steps early in the run.
    Step counts are comparable only within a regime and within a chip; hatched bars are chip~E, which ran slow.}
  \label{fig:steps}
\end{figure}

\clearpage
\section{Oracle details}
\label{app:oracle}

\begin{table}[!htbp]
  \centering
  \caption{\textbf{Offset $=$ official $-$ rehearsal, by period} ($10^{-3}$ bpb). The 95\% CI is the half-width for
    the mean (Student $t$). Ranges are rounded half-up from the exact differences. Official scores are published to
    four decimals, so every offset carries a rounding error of up to $\pm$\val{official-rounding}.}
  \label{tab:offset-periods}
  % GENERATED by paper/analysis -- offset = official - rehearsal by period, 1e-3 bpb (oracle.py)
{\small
\begin{tabular}{@{}p{5.6cm}rrrrr@{}}
\toprule
Uploads & $n$ & Mean & SD & Range & 95\% CI \\
\midrule
chip-C calibration uploads, K51--K60 & 7 & \ensuremath{+}6.53 & 0.42 & \ensuremath{+}6.1 to \ensuremath{+}7.2 & $\pm$0.39 \\
final day, chips C/E/G/H, K63--K82s7 & 7 & \ensuremath{+}6.87 & 0.30 & \ensuremath{+}6.5 to \ensuremath{+}7.4 & $\pm$0.27 \\
since K50 (all chips) & 15 & \ensuremath{+}6.71 & 0.38 & \ensuremath{+}6.1 to \ensuremath{+}7.4 & $\pm$0.21 \\
early, chips A/B, F2--K46 & 15 & \ensuremath{+}6.26 & 0.72 & \ensuremath{+}5.3 to \ensuremath{+}7.8 & $\pm$0.40 \\
every rehearsed upload & 30 & \ensuremath{+}6.48 & 0.61 & \ensuremath{+}5.3 to \ensuremath{+}7.8 & $\pm$0.23 \\
\bottomrule
\end{tabular}
}

\end{table}

\begin{table}[!htbp]
  \centering
  \caption{\textbf{Coverage of the expanding-window prediction intervals}: uploads inside the interval, with the
    expected count in brackets, for two burn-in lengths (\cref{sec:oracle-holdout}).}
  \label{tab:coverage}
  % GENERATED by paper/analysis -- observed (expected) count of uploads inside the expanding-window interval (oracle.py)
{\small
\begin{tabular}{@{}lcc@{}}
\toprule
Nominal level & Burn-in 3 uploads & Burn-in 5 uploads \\
\midrule
50\% & 7/12 (6.0) & 6/10 (5.0) \\
80\% & 11/12 (9.6) & 9/10 (8.0) \\
95\% & 12/12 (11.4) & 10/10 (9.5) \\
\bottomrule
\end{tabular}
}

\end{table}

\begin{table}[!htbp]
  \centering
  \caption{\textbf{Sensitivity of the oracle statistics to \upl{K70a}} ($10^{-3}$ bpb). ``Frozen'': the chip-C
    calibration on the final-day uploads. ``In 95\% PI'' and ``Seq.\ MAE'': the expanding window with a burn-in of
    three uploads.}
  \label{tab:oracle-sens}
  % GENERATED by paper/analysis -- oracle statistics with K70a as used, raw and excluded, 1e-3 bpb (oracle.py)
{\footnotesize
\begin{tabular}{@{}p{4.9cm}rrrrrrr@{}}
\toprule
Data since K50 & $n$ & Mean & SD & Frozen MAE & Frozen bias & In 95\% PI & Seq.\ MAE \\
\midrule
K70a normalised to chip-C speed (as used) & 15 & \ensuremath{+}6.71 & 0.38 & 0.35 & \ensuremath{+}0.34 & 12/12 & 0.31 \\
K70a raw chip-E rehearsal & 15 & \ensuremath{+}6.62 & 0.54 & 0.53 & \ensuremath{+}0.15 & 11/12 & 0.47 \\
K70a excluded & 14 & \ensuremath{+}6.72 & 0.39 & 0.40 & \ensuremath{+}0.39 & 11/11 & 0.34 \\
\bottomrule
\end{tabular}
}

\end{table}

\subsection{What the offset variance is made of}
\label{app:offset-detail}

Since \upl{K50} the within-chip offset SD is \val{off-sd-within}. If the official text behaves like the first 20M public tokens, the run-to-run variation of
the text gap (SD \val{offdec-text}) explains only \val{offdec-text-share} of the within-chip offset variance and
rounding to four decimals almost none (SD \val{offdec-round}); the remaining SD, \val{offdec-rest}
(\val{offdec-rest-lo}--\val{offdec-rest-hi} at the limits of the within-chip SD's interval), comes from step-count
jitter, from the official side or from nondeterminism, which these data cannot separate.

\paragraph{Chip or model quality?} On the final day the faster chips also rehearsed the better recipes and the
lottery winners, so an apparent chip effect could be an effect of model quality or of selection. Across the
\val{cq-n} uploads since \upl{K50}, the offset's slope on the rehearsal score is \val{cq-slope} (95\% CI
\val{cq-slope-lo} to \val{cq-slope-hi}, $p=\val{cq-slope-p}$), which moves the offset by \val{cq-slope-span} across
the rehearsal span. The \val{cq-gh-n} uploads rehearsed on chips G and H have a mean offset \val{cq-gh-diff} above the
rest (exact permutation $p=\val{cq-perm-p}$ over all \val{cq-perm-n} label assignments), but given the rehearsal score
the G/H term is \val{cq-gh-coef} (95\% CI \val{cq-gh-lo} to \val{cq-gh-hi}; permutation $p=\val{cq-perm-p-adj}$) and
the slope becomes \val{cq-slope-adj}. With so few uploads the data cannot tell a chip effect from an effect of model
quality or selection, so we say only that the offset \emph{may} depend on the chip.

\subsection{Traps in detail}
\label{app:traps}

The four ways in which a rehearsal was not a faithful replica of the official run (\cref{sec:traps}):
\begin{enumerate}
  \item \textbf{Cold means cold.} Any one-time compilation that happens after the clock starts costs steps in the
    official run. A warm rehearsal, which reuses compiled graphs, does not pay that cost. \Cref{sec:prewarm}
    quantifies one instance: warm rehearsals of an EMA recipe looked \val{ema-warm-bias}\docnum{} bpb better than the
    cold official run would be.
  \item \textbf{Chip speed.} A rehearsal predicts the official score only if the official host runs at the
    rehearsal chip's speed; otherwise $\gamma_c\neq0$. We normalised same-recipe runs to chip-C speed at $\rho$ and
    never across a change that alters step time.
  \item \textbf{The warm-up coin flip.} A change that touches steps 0--19 can move a seed onto a different early
    trajectory (\cref{sec:seeds}). The rehearsal still predicts the official score, because the official run takes
    the same path, but the measured difference is then a seed effect rather than the change's effect.
  \item \textbf{The time target is a budget, not a promise.} The training clock stops when the next step would
    cross the target; the median reserve was \val{reserve-s}\docnum{}~s over \val{reserve-n} completed runs, and we
    kept \val{headroom}\docnum{}~s of headroom under the cap.
\end{enumerate}
Two scores from runs that the organisers rescored, one each for \upl{K44} and \upl{F6}, are excluded from all
calibration statistics, so that each upload counts once; both are listed in
\path{research/sim-data/official-uploads.csv}.

\subsection{The normalised rehearsal of \upl{K70a}}
\label{sec:k70a}

Chip~E ran about \val{chip-spread}\docnum{} slower in step time than chip~C. \upl{K70a}'s raw chip-E rehearsal,
\val{k70a-raw}, was converted to chip-C speed by subtracting $\rho$ times its step deficit, which gave
\val{k70a-norm}: an adjustment of \val{k70a-adj}, i.e.\ \val{k70a-adj-pct} of steps at $\rho$ (not all of a
step-time deficit becomes a step deficit, because the schedule is clock-keyed). Its raw offset would be
\val{k70a-raw-offset}, an outlier. The repository does not show whether the normalised value was computed before the
upload: no time-stamped projection exists, so we treat it as a post-hoc adjustment. \Cref{tab:oracle-sens} gives
every oracle statistic three ways. With the raw rehearsal the since-\upl{K50} SD rises from \val{sens-used-sd} to
\val{sens-raw-sd}, the frozen MAE from \val{sens-used-mae} to \val{sens-raw-mae}, and the expanding-window count
falls to \val{sens-raw-inside}; excluding \upl{K70a} changes little.

\subsection{Expanding-window prediction}

\Cref{fig:offset-prospective} predicts each upload from the mean offset of all earlier uploads since \upl{K50}, with
a $t$ prediction interval, starting once three earlier points exist (\cref{sec:oracle-holdout}). \Cref{tab:coverage}
reports coverage at 50, 80 and 95\% for burn-ins of three and five uploads: at 50\% the intervals cover
\val{cov-50-b3} and \val{cov-50-b5}, close to nominal, and at 80\% \val{cov-80-b3} and \val{cov-80-b5}, slightly
conservative.

\begin{figure}[!htbp]
  \centering
  \includegraphics{offset_prospective.pdf}
  \caption{\textbf{Pseudo-prospective offset prediction.} Each point is an upload's offset (marker: rehearsal chip;
    grey: the three burn-in uploads). The solid line is the mean offset of all earlier uploads since \upl{K50}, the
    band its 95\% prediction interval, and the dashed line the frozen chip-C fit. Generated by
    \texttt{paper/analysis/oracle.py}.}
  \label{fig:offset-prospective}
\end{figure}

\subsection{The written projections}
\label{sec:projections}

The repository records three projections written before upload (\texttt{results/official-scores.csv},
\texttt{docs/EXACT-ORACLE.md}). \upl{K73s6} was projected at \val{k73s6-projected}\docnum{} and scored
\val{official-k73s6}, the largest miss (\val{k73s6-miss}). \upl{K82s4} was projected at
\val{k82s4-projected}\docnum{} and scored \val{best-official}; \upl{K82s7} was projected at about \val{k82s7-projected}\docnum{} and
scored \val{salt7-official}. These two projections did \emph{not} use the frozen chip-C offset, which misses them by
\val{ferr-k82s4} and \val{ferr-k82s7} (\cref{tab:offset-holdout}). They added an offset of \val{k82-proj-offset}\docnum{}
read off the final-day uploads already scored, most of them on the faster chips G and H; the exact rule was not
recorded. That is the text-plus-chip correction applied informally, which may be why those two landed close. It is
also a reason to treat their closeness as anecdote rather than as a test.

\section{Simulator internals}
\label{app:sim}

\subsection{Dataset}
\label{app:simdata}

Each record of \texttt{runs.jsonl} carries the effective \texttt{FF\_*} environment, steps, charged and start-up
seconds, median step time per accumulation phase, the early loss curve where logged (\val{ds-loss-curves} records),
and the \valbpb{} on the first 2M public tokens. Records outside the quality fit are screens, runs from the earliest
chips without a recorded code version, rows whose base recipe is unknown, or runs without a final score. The pairs
of \texttt{validation-pairs.json} pair each treatment with the control the campaign named at the time, leaving out
60-step screens\docnum{}; no earlier comparisons were paired.

\begin{table}[!htbp]
  \centering
  \caption{\textbf{Run records by instance} (\texttt{research/sim-data/runs.jsonl}). Chips A and B are the early
    instances whose records come from a run table without a code version; chip D ran an older Neuron runtime
    (2.33.10, against 2.34.10 on chip C). ``Unknown'': records without a recorded chip. The final-day chips E, G and
    H are not in the dataset.}
  \label{tab:dataset}
  % GENERATED by paper/analysis -- runs.jsonl by chip (simulator.py)
{\small
\begin{tabular}{@{}lrrrrr@{}}
\toprule
Chip & Records & Full runs & Screens & Full, with bpb & Loss curves \\
\midrule
unknown & 28 & 28 & 0 & 5 & 0 \\
A & 563 & 426 & 137 & 390 & 0 \\
B & 247 & 234 & 13 & 220 & 0 \\
C & 349 & 223 & 126 & 215 & 331 \\
D & 19 & 17 & 2 & 17 & 11 \\
\midrule
all & 1{,}206 & 928 & 278 & 847 & 342 \\
\bottomrule
\end{tabular}
}

\end{table}

\subsection{Step time and the step walk}
\label{app:steptime}

The step-time model regresses the median $k\!=\!1,2,4$ step times on code-lineage effects and knob features (depth,
width, batch, fused paths). It is deliberately empirical, because the Neuron compiler is far from linear in
operation count: removing one elementwise pass over the MLP activation saved \val{st-onepass}\docnum{}, fusing the
QKV projection, which has fewer operations, was \val{st-qkv}\docnum{} \emph{slower}, and a 96-parameter gate cost
\val{st-gate}\docnum{}. A mechanism that no fitted run carries is charged a \val{st-novel-prior}\docnum{} prior and
flagged ``cost unknown''. Medians are scaled by the measured median-to-mean overhead (\val{overhead}) before the walk.
With one code lineage left out, the model predicts the $k\!=\!4$ step time of runs whose mechanisms another lineage
carried, with an MAE of \val{st-lolo-reach} (\val{st-lolo-nreach} runs; \cref{tab:sim-other}), and the walk reproduces
a rehearsal's phase switches to within \val{walk-phase-err}\docnum{} steps.

\subsection{Features and code eras}

The quality surrogate reads \val{q-nbase-features} base features (\cref{tab:features}) and \val{q-nbowl-features}
non-negative curvature (``bowl'') terms for tuned scalars, the squares of the corresponding base features. Every
feature is centred on the \upl{K59} reference recipe, so the era intercept of the current code family is that recipe's
predicted score at 2{,}300 steps on chip~C. Each feature's scale in \cref{tab:features} is the size of a typical campaign move; the ridge prior applies
per scaled unit. \emph{Code eras} are groups of versions of our training script: M3 to M14 number the main lineage in
order, and ``M9+'' groups the flag-gated versions from M9 on, whose new mechanisms are off unless a knob turns them on.

\begin{table}[!htbp]
  \centering
  \caption{\textbf{Knob features of the quality surrogate} (\texttt{ffsim/quality.py}, \texttt{FEATURES}). Two
encodings were set with campaign outcomes in view: the \texttt{leaky\_sq} bowl, centred on the slope 0.35 that had
measured best, and \texttt{acc\_in\_graph}, added after the \upl{K56} mechanism had measured well.}
  \label{tab:features}
  % GENERATED by paper/analysis -- ffsim.quality FEATURES (simulator.py)
{\scriptsize
\begin{tabular}{@{}lp{10.6cm}r@{}}
\toprule
Feature & Encoding (centred on the K59 recipe) & Scale \\
\midrule
\texttt{leaky} & LK - 0.35 (leaky relu\textsuperscript{2} slope; 0 = plain relu\textsuperscript{2}) & 0.15 \\
\texttt{leaky\_sq} & (LK - 0.35)\textsuperscript{2}: a bowl in the slope (0 / 0.25 / 0.35 / 0.5 were tried) & 0.0225 \\
\texttt{leaky\_form\_abs} & 1 if LK > 0 and form is 'abs' (M10b); fn/fnm = 0 & 1 \\
\texttt{leaky\_form\_max} & 1 if LK > 0 and form is 'max' (M11c, torch.maximum) & 1 \\
\texttt{relu2\_fn} & 1 if FF\_RELU2\_FN (the RF arm: relu\textsuperscript{2} with a hand-written backward, LK = 0) & 1 \\
\texttt{softcap\_log} & ln(SOFTCAP / 15) when SOFTCAP > 0, else 0 & 0.182 \\
\texttt{softcap\_off} & 1 if SOFTCAP == 0 (no logit cap) & 1 \\
\texttt{softcap\_a} & A\_eff / SOFTCAP - 1 with A\_eff = SOFTCAP\_A if > 0 else SOFTCAP (asymmetric cap A*tanh(z/C): 16.5 on 15 $\rightarrow$ 0.1) & 0.1 \\
\texttt{softcap\_b} & SOFTCAP\_B (logit shift inside the tanh) & 5 \\
\texttt{wd\_sched} & WD\_SCHED power (muon wd = WD * lrm ** WD\_SCHED; 0 = constant wd; non-numeric $\rightarrow$ 1) & 1 \\
\texttt{mom\_peak} & MOM\_PEAK - 0.95 (Muon momentum peak; 0.93 / 0.97 were tried) & 0.02 \\
\texttt{muon\_beta2} & MUON\_BETA2 - 0.9 (0.8 and 0.98 were tried) & 0.05 \\
\texttt{rope\_log} & ln(ROPE\_BASE / 1e5) (300k tried) & 1.1 \\
\texttt{scalar\_beta2} & SCALAR\_BETA2 - 0.95 (0.99 tried) & 0.04 \\
\texttt{qk\_gain} & QK\_GAIN - 1.44 (1.6 tried) & 0.16 \\
\texttt{warmup\_log} & ln(WARMUP / 20) (40 tried) & 0.693 \\
\texttt{emb\_lr\_log} & ln(EMB\_LR / 0.30) (0.36 tried) & 0.182 \\
\texttt{unemb\_lr\_log} & ln(UNEMB\_LR / 0.006) (0.0072 tried) & 0.182 \\
\texttt{scalar\_lr\_log} & ln(SCALAR\_LR / 0.20) (0.1664 tried with ADAMW\_LR 1.7) & 0.182 \\
\texttt{matrix\_lr\_log} & ln(MATRIX\_LR\_SCALE / 2.0) (1.75 / 2.3 tried) & 0.14 \\
\texttt{adamw\_lr\_log} & ln(ADAMW\_LR\_SCALE / 1.4142) (1.7 tried) & 0.182 \\
\texttt{wd\_log} & ln(WD / 0.02) (0.012 / 0.03 tried) & 0.405 \\
\texttt{key\_offset\_off} & 1 if the rotary key offset is off (pre-K51 recipes) & 1 \\
\texttt{key\_offset\_log} & ln(radius / 0.3) when the key offset is on (0.03 / 0.1 / 0.2 / 0.5 / 1.0 / 3.0 were tried: two decades, so a log scale), 0 when off & 1.1 \\
\texttt{acc\_in\_graph} & ACC\_IN\_GRAPH - 1 (K56 mechanism: gradient accumulation inside the compiled graph) & 1 \\
\texttt{accum\_until} & charged-time fraction at which the batch warm-up reaches the final k, minus 0.2 (plain 'K' $\rightarrow$ -0.2; '2:0.2,4' and RAMP3 $\rightarrow$ 0) & 0.2 \\
\texttt{accum\_k1\_until} & first-phase (k1) end fraction minus 0.06 (RAMP3 $\rightarrow$ 0; two-phase or none $\rightarrow$ -0.06) & 0.06 \\
\texttt{accum\_lr0} & LR scale of the first small-batch phase minus 0.5 (0 when no batch warm-up) & 0.25 \\
\texttt{accum\_lr1} & LR scale of the second small-batch phase minus 0.7071 (0 unless RAMP3) & 0.25 \\
\texttt{cooldown\_frac} & COOLDOWN\_FRAC - 0.6 (0.45 / 0.55 / 0.65 tried) & 0.1 \\
\texttt{mtp\_off} & 1 if FF\_MTP == 0 (no multi-token head) & 1 \\
\texttt{mtp\_p1} & MTP fade phase 1 - 0.2 (0 when MTP off) & 0.1 \\
\texttt{mtp\_p2} & MTP fade phase 2 - 0.45 (0 when MTP off; 0.5 tried) & 0.1 \\
\texttt{mtp\_w\_tail} & sum of the MTP head weights after the first, minus 0.75 ('1,0.6,0.3' $\rightarrow$ 0.15; 0 when MTP off) & 0.15 \\
\texttt{xsa\_on} & 1 if FF\_XSA\_LAYERS is non-empty & 1 \\
\texttt{xsa\_frac} & fraction of blocks with XSA ('all' $\rightarrow$ 1; '0,1,2,3' on d9 $\rightarrow$ 0.44) & 1 \\
\texttt{attn\_src} & 1 if FF\_ATTN\_SRC is non-empty (attention-source reuse, M14 '5:6,7,8') & 1 \\
\texttt{fn\_flags} & fraction of the M13 autograd.Function flags on (step-time mechanisms, quality-neutral by design: prior sd 0.001) & 1 \\
\texttt{ce\_bf16} & 1 if the cross-entropy runs in bf16 (G3 probe, rejected) & 1 \\
\texttt{ptp\_w} & PTP\_W (previous-token-prediction head weight; 0.25 tried) & 0.25 \\
\texttt{mom\_ramp\_log} & ln(MOM\_RAMP / 300) (450 tried) & 0.405 \\
\texttt{depth} & DEPTH - 9 (8 / 10 / 12 tried) & 1 \\
\texttt{aspect\_log} & ln(ASPECT\_RATIO / 113) (width = aspect * depth) & 0.0953 \\
\texttt{total\_batch\_log2} & log2(TOTAL\_BATCH / 262144) & 1 \\
\texttt{mlp\_mult} & MLP\_MULT - 4 & 1 \\
\texttt{time\_target\_log} & ln(TIME\_TARGET / 1793): mostly absorbed by ln(steps), so its prior sd is 0.001 and it only catches what the step count misses & 0.0953 \\
\bottomrule
\end{tabular}
}

\end{table}

\subsection{Fitting the quality surrogate}
\label{app:fit}

The fit of \cref{eq:quality} is not one closed-form solve. For a given residual SD $\sigma_q$, the MAP estimate solves
$(X^\top X + \sigma_q^2 D)\hat\beta = X^\top y + \sigma_q^2 D m$ with $D=\mathrm{diag}(1/s_j^2)$ for prior SDs $s_j$
and prior means $m$. A bowl coefficient that comes out negative is pinned at zero and the system re-solved (an
active-set solution of the non-negativity constraint). $\sigma_q$ is then re-estimated by empirical Bayes from the
residuals and effective degrees of freedom, three times, and floored at \val{q-sigma-floor}. Seed effects are
estimated jointly, as best linear unbiased predictors. The fit uses \val{q-nfit} runs, \val{q-ncols} columns and
\val{q-neras} eras, with \val{q-edf} effective degrees of freedom.

\paragraph{Paired Monte Carlo.} \texttt{simulate} and \texttt{search} propagate chip jitter, seed luck, coefficient
uncertainty, residuals and the offset by Monte Carlo with common random numbers, so that wherever two recipes agree
the noise cancels. $P(\text{candidate} < \text{base})$ is a paired probability; we have not checked its calibration
beyond interval coverage (\val{lopo-rec-cover} of LOPO recipe pairs fall inside the predicted $\pm2$\,SD,
\cref{sec:simval}). An identical recipe reproduces the base exactly.

\subsection{What the surrogate learned}

\Cref{tab:coefficients} lists the knob coefficients furthest from zero in posterior-SD units. Read it with the
support columns. A feature varied in a single run (\texttt{ce\_bf16}, \texttt{accum\_lr1}) or only inside early eras
(\texttt{key\_offset\_off}, varied only in M3 and M4) is identified mostly by those runs and partly absorbs era
differences. Such coefficients are bookkeeping for prediction, not causal effect sizes. The attention-source
coefficient (\texttt{attn\_src}) rests on two runs, which is why the post-fit attention-source pairs were
over-predicted (\cref{tab:postfit}).

\begin{table}[!htbp]
  \centering
  \caption{\textbf{The \val{coef-ntop} knob coefficients furthest from zero} in posterior-SD units, in $10^{-3}$ bpb
    per scaled unit. ``Runs $\neq0$'' counts fitted runs in which the feature is non-zero, and ``Code eras'' the eras
    those runs come from. Knob names omit the \texttt{FF\_} prefix.}
  \label{tab:coefficients}
  % GENERATED by paper/analysis -- the 16 knob coefficients furthest from zero in posterior-SD units, x 1e-3 bpb per scaled unit
{\footnotesize\setlength{\tabcolsep}{4pt}
\begin{tabular}{@{}lp{3.6cm}rrrrp{2.6cm}@{}}
\toprule
Feature & Knob(s) & Est. & Post.\ SD & Prior SD & Runs $\neq0$ & Code eras \\
\midrule
\texttt{accum\_until} & ACCUM\_SCHED & \ensuremath{+}6.37 & 0.60 & 4.0 & 10 & M3, M6, M7 \\
\texttt{ce\_bf16} & CE\_BF16 & \ensuremath{+}2.93 & 0.34 & 4.0 & 1 & M6 \\
\texttt{attn\_src} & ATTN\_SRC & \ensuremath{-}2.02 & 0.25 & 4.0 & 2 & M9+ \\
\texttt{accum\_lr0} & ACCUM\_LR, ACCUM\_SCHED & \ensuremath{-}2.42 & 0.44 & 4.0 & 34 & M3, M4, M5, M6, M7 \\
\texttt{key\_offset\_log\_sq} & KEY\_OFFSET\_RAD, KEY\_OFFSET & \ensuremath{+}0.46 & 0.10 & 1.0 & 48 & M4, M5, M6, M7 \\
\texttt{leaky\_sq} & LEAKY\_RELU2 & \ensuremath{+}0.28 & 0.06 & 1.0 & 111 & M3, M4, M5, M6, M7, M9+ \\
\texttt{mom\_peak\_sq} & MOM\_PEAK & \ensuremath{+}1.09 & 0.24 & 1.0 & 2 & M9+ \\
\texttt{mom\_peak} & MOM\_PEAK & \ensuremath{+}0.74 & 0.21 & 4.0 & 2 & M9+ \\
\texttt{key\_offset\_off} & KEY\_OFFSET & \ensuremath{-}4.75 & 1.41 & 4.0 & 8 & M3, M4 \\
\texttt{accum\_k1\_until} & ACCUM\_SCHED & \ensuremath{+}0.90 & 0.34 & 4.0 & 35 & M3, M4, M5, M6, M7 \\
\texttt{accum\_lr1} & ACCUM\_LR, ACCUM\_SCHED & \ensuremath{+}2.43 & 0.94 & 4.0 & 1 & M7 \\
\texttt{xsa\_frac} & XSA\_LAYERS, DEPTH & \ensuremath{-}1.84 & 0.86 & 4.0 & 2 & M9+ \\
\texttt{leaky} & LEAKY\_RELU2 & \ensuremath{+}0.18 & 0.11 & 4.0 & 111 & M3, M4, M5, M6, M7, M9+ \\
\texttt{softcap\_a} & SOFTCAP\_A, SOFTCAP & \ensuremath{-}0.29 & 0.20 & 4.0 & 4 & M5, M9+ \\
\texttt{mom\_ramp\_log} & MOM\_RAMP & \ensuremath{+}0.35 & 0.42 & 4.0 & 2 & M5 \\
\texttt{matrix\_lr\_log} & MATRIX\_LR\_SCALE & \ensuremath{+}0.18 & 0.22 & 4.0 & 2 & M7 \\
\bottomrule
\end{tabular}
}

\end{table}

\subsection{Pair validation by family, forward in time, and on new mechanisms}
\label{app:pairval}

Against the family-mean baseline, which scores \val{base-fam-rec-sign} (\val{base-fam-rec-signci}) with MAE
\val{base-fam-rec-mae} on the recipe pairs, the surrogate's MAE difference, \val{base-diff-fam-mae-ci}$\times10^{-3}$,
includes zero. ``Decisive'' pairs defined by the measured outcome ($|\Delta|\ge0.0003$) favour large step-count changes, so
\cref{sec:simval} reports the pairs on which the \emph{prediction} is decisive; on them the MAE is
\val{base-full-pdec-mae} for the surrogate against \val{base-steps-pdec-mae} for the steps-only model. Forward in
time, the surrogate's MAE is better than the steps-only model's with the 29~September cutoff (difference
\val{temp-b-diff-mae-ci}$\times10^{-3}$) but not with 28~September (\val{temp-a-diff-mae-ci}$\times10^{-3}$), where the zero
predictor's MAE is \val{temp-a-zero-mae}. During development (\cref{tab:surrogate-configs}), a fix to how missing
knobs are read moved LOPO sign agreement on the pairs named in the simulator's specification from
\val{sc-spec-sign-before}\docnum{} to \val{sc-spec-sign-after}\docnum{}, and recipe-pair sign agreement ranged from
\val{sc-rec-sign-lo}\docnum{} (\val{sc-rec-lo-name}) to \val{sc-rec-sign-hi}\docnum{} (\val{sc-rec-hi-name}) across
configurations; the default, \val{sc-default-name}, scored \val{sc-default-rec-sign}\docnum{}.

\begin{table}[!htbp]
  \centering
  \caption{\textbf{Leave one pair out (LOPO) and leave one knob value out (LOKVO) by pair family}, in $10^{-3}$ bpb.
    The ``seed'' and ``chip'' families compare identical recipes and measure noise; ``cold upload rehearsals'' pairs
    compare upload rehearsals.}
  \label{tab:lopo-family}
  % GENERATED by paper/analysis -- per-family pair validation, 1e-3 bpb (simulator.py)
{\footnotesize
\begin{tabular}{@{}p{4.4cm}rrrrrr@{}}
\toprule
Pair family & $n$ & Mean $\Delta$ & LOPO sign & LOPO MAE & LOKVO sign & LOKVO MAE \\
\midrule
LeakyReLU slope & 13 & \ensuremath{-}0.83 & 92\% & 0.33 & 69\% & 0.82 \\
seed (noise) & 11 & \ensuremath{+}1.43 & 100\% & 0.59 & 100\% & 0.59 \\
scalar retunes, K57 base & 10 & \ensuremath{+}0.28 & 70\% & 0.71 & 70\% & 0.73 \\
batch warm-up schedule & 9 & \ensuremath{-}0.23 & 89\% & 0.17 & 89\% & 0.23 \\
cold upload rehearsals & 9 & \ensuremath{-}0.64 & 67\% & 0.27 & 78\% & 0.96 \\
LeakyReLU form & 8 & \ensuremath{+}0.03 & 88\% & 0.41 & 88\% & 0.99 \\
key-offset radius & 7 & \ensuremath{-}0.20 & 100\% & 0.20 & 100\% & 0.19 \\
weight-decay schedule & 4 & \ensuremath{-}0.03 & 75\% & 0.08 & 75\% & 0.07 \\
accumulation in graph & 3 & \ensuremath{-}0.61 & 100\% & 0.06 & 100\% & 0.08 \\
QK gain & 3 & \ensuremath{+}0.08 & 33\% & 0.59 & 33\% & 0.54 \\
ReLU$^2$ hand-written backward & 3 & \ensuremath{-}1.12 & 100\% & 0.16 & 100\% & 0.10 \\
AdamW $\beta_2$ & 2 & \ensuremath{+}0.39 & 100\% & 0.96 & 100\% & 0.96 \\
learning rate & 2 & \ensuremath{+}0.18 & 100\% & 0.45 & 100\% & 0.45 \\
Muon momentum peak & 2 & \ensuremath{+}1.26 & 50\% & 1.21 & 50\% & 1.21 \\
asymmetric soft cap & 2 & \ensuremath{-}0.08 & 50\% & 0.17 & 100\% & 0.11 \\
attention-input shift & 2 & \ensuremath{+}2.30 & 100\% & 0.77 & 100\% & 0.77 \\
ZeRO-2 + AdamW ZeRO & 2 & \ensuremath{-}1.39 & 100\% & 0.19 & 100\% & 0.19 \\
chip (noise) & 1 & \ensuremath{+}4.68 & 100\% & 0.39 & 100\% & 0.39 \\
code-M13 flags & 1 & \ensuremath{+}0.06 & 0\% & 0.32 & 0\% & 0.32 \\
attention-source reuse & 1 & \ensuremath{-}1.73 & 100\% & 0.42 & 100\% & 1.68 \\
multi-token prediction & 1 & \ensuremath{+}0.06 & 100\% & 0.04 & 100\% & 0.04 \\
RoPE base & 1 & \ensuremath{+}0.15 & 100\% & 0.17 & 100\% & 0.17 \\
soft cap & 1 & \ensuremath{+}0.35 & 100\% & 0.08 & 100\% & 0.08 \\
LR warm-up & 1 & \ensuremath{+}0.66 & 100\% & 0.68 & 100\% & 0.68 \\
\bottomrule
\end{tabular}
}

\end{table}

\begin{figure}[!htbp]
  \centering
  \includegraphics{simulator_pairs.pdf}
  \caption{\textbf{The quality surrogate on pairs it did not see} (leave one pair out, all \val{lopo-all-n} pairs;
    values beyond $\pm4\times10^{-3}$ are drawn at the edge). Shaded quadrants mark a wrong predicted sign. Squares
    are seed and chip pairs, which measure run-to-run noise rather than a recipe effect.}
  \label{fig:pairs}
\end{figure}

\begin{table}[!htbp]
  \centering
  \caption{\textbf{Temporal hold-out of the quality surrogate.} ``Supported'' pairs move only features that some run
    in the restricted fit varied; ``Prior-only'' pairs move a feature that no such run varied, so the model can only
    return its prior for it. Cells in these two columns give $n$: sign agreement of the surrogate. ``Steps only'': the same model
    with every knob effect switched off, fitted on the same early runs; ``Zero'': the predictor $\hat\Delta=0$.}
  \label{tab:temporal}
  % GENERATED by paper/analysis -- temporal hold-out (simulator.py); MAE in 1e-3 bpb; supported / prior-only cells are n: sign agreement
{\footnotesize\setlength{\tabcolsep}{4pt}
\begin{tabular}{@{}lrrrrrrrrr@{}}
\toprule
MAE in $10^{-3}$ bpb & & & \multicolumn{2}{c}{Surrogate} & \multicolumn{2}{c}{Steps only} & Zero & & \\
\cmidrule(lr){4-5}\cmidrule(lr){6-7}
Fit on & Runs & Pairs & Sign & MAE & Sign & MAE & MAE & Supported & Prior-only \\
\midrule
runs before 28 Sep & 59 & 72 & 86\% & 0.89 & 71\% & 0.78 & 0.94 & 36: 89\% & 36: 83\% \\
runs before 29 Sep & 98 & 46 & 91\% & 0.48 & 67\% & 0.89 & 1.08 & 32: 91\% & 14: 93\% \\
\bottomrule
\end{tabular}
}

\end{table}

With the 28~September cutoff, the \val{temp-a-sup-n} temporal pairs whose moved features the restricted fit had
seen score \val{temp-a-sup-sign}, and the \val{temp-a-uns-n} pairs that only the prior can answer score
\val{temp-a-uns-sign}; with 29~September the two groups score \val{temp-b-sup-sign} and \val{temp-b-uns-sign}
(\cref{tab:temporal}). For the prior-only pairs the sign must come from the step term, the era and seed effects and
the positive prior means of the bowl terms, which make a move away from a tuned value cost (\cref{app:fit}).

\begin{table}[!htbp]
  \centering
  \caption{\textbf{The \val{postfit-n} pairs that finished after the release fit} ($10^{-3}$ bpb). Observed and
    predicted values per pair are from \texttt{docs/simulator-report-20260929.md} \S5\docnum{} (the model was the
    release fit, in the candidate convention); sign agreement, the Wilson interval and the MAE are computed by
    \texttt{paper/analysis/simulator.py}; the Wilson interval treats the pairs as independent, but several share a
    control, so it is too narrow. AS: attention-source variants; ALR: AdamW learning-rate change.}
  \label{tab:postfit}
  % GENERATED by paper/analysis -- 15 post-fit pairs, 1e-3 bpb (simulator.py)
{\footnotesize
\begin{tabular}{@{}p{6.4cm}rrc@{}}
\toprule
Pair (report label) & Observed & Predicted & Sign \\
\midrule
C41 K60 vs C40 K59 (cold rehearsals) & \ensuremath{-}1.82 & \ensuremath{-}2.10 & yes \\
AS-a s67 (C) & \ensuremath{-}1.36 & \ensuremath{-}2.04 & yes \\
AS-a s58 (C) & \ensuremath{-}1.06 & \ensuremath{-}1.71 & yes \\
AS-a s97 (D) & \ensuremath{+}1.15 & \ensuremath{-}0.99 & no \\
AS-b 4:5,6,7,8 s73 (D) & \ensuremath{-}0.86 & \ensuremath{-}2.29 & yes \\
K60 + ALR s73 (C) vs AS-a & \ensuremath{+}0.81 & \ensuremath{+}1.67 & yes \\
K60 + ALR s67 (C) vs AS-a & \ensuremath{+}0.95 & \ensuremath{+}2.24 & yes \\
ALR s97 (D) vs LK0.35 & \ensuremath{-}0.78 & \ensuremath{+}1.73 & no \\
ALR s113 (D) vs LK0.35 & \ensuremath{-}0.09 & \ensuremath{+}1.81 & no \\
ALR s131 (D) vs LK0.35 & \ensuremath{+}0.44 & \ensuremath{+}2.13 & yes \\
AS-d 5:7,8 s73 (C) vs AS-a & 0.00 & \ensuremath{+}0.15 & no \\
AS-d 5:7,8 s67 (C) vs AS-a & \ensuremath{+}0.33 & \ensuremath{+}0.39 & yes \\
AS-d 5:7,8 s73 (D) vs AS-a & \ensuremath{-}0.20 & \ensuremath{+}0.02 & no \\
AS-c 6:7,8 s73 (D) vs AS-a & \ensuremath{+}0.76 & \ensuremath{-}0.07 & no \\
K60 + softcap 16.5 s73 (C) vs AS-a & \ensuremath{+}0.68 & \ensuremath{-}0.22 & no \\
\midrule
sign agreement 8/15 (Wilson 95\%: 30\%--75\%) & & MAE 1.04 & \\
\bottomrule
\end{tabular}
}

\end{table}

\begin{table}[!htbp]
  \centering
  \caption{\textbf{Step-time model, end-to-end anchors and the offset model.} End-to-end rows compare the Monte Carlo
    prediction at seed~73 with the measured upload (\texttt{docs/SIMULATOR.md}\docnum{}).}
  \label{tab:sim-other}
  % GENERATED by paper/analysis -- step time, end-to-end anchors and the offset model (simulator.py); anchors from docs/SIMULATOR.md
{\footnotesize
\begin{tabular}{@{}p{6.2cm}rll@{}}
\toprule
Check & Data & Predicted & Measured / note \\
\midrule
Step time, k4 phase, leave one code lineage out: runs whose mechanisms another lineage ran & 161 runs & MAE 0.54\% & target 0.5\% \\
\quad all held-out runs, including novel mechanisms & 189 runs & MAE 0.95\% & up to 14\% \\
\midrule
End to end, \upl{K59}: steps / rehearsal / official$^{\dagger}$ & 1 upload & 2{,}369 / 0.96082 / 0.9674 & 2{,}357 / 0.96092 / 0.9671 \\
End to end, \upl{K60}: steps / rehearsal / official$^{\dagger}$ & 1 upload & 2{,}365 / 0.95888 / 0.9654 & 2{,}361 / 0.95910 / 0.9655 \\
\midrule
Offset model (chip-C uploads): mean and SD & 7 uploads & \ensuremath{+}0.00653, 0.00042 & SD 95\% CI 0.00027--0.00092 \\
\bottomrule
\end{tabular}
}

\end{table}

The full validation reports, including support grading for every knob value and the leakage guard, are
\texttt{research/sim-data/steptime-validation.md}, \texttt{simulate-validation.md} and \texttt{dataset-qa.md}.
\texttt{research/sim-data/quality-validation.md} is a historical report: it was written on an earlier snapshot of the
dataset (\val{qv-records}\docnum{} records), and the script that generated it is not released, so its numbers are
not recomputed here. Its \S11 compares the \val{sc-n}\docnum{} surrogate configurations of \cref{tab:surrogate-configs},
whose LOPO results were visible while the default was chosen (\cref{sec:simval}).

\begin{table}[!htbp]
  \centering
  \caption{\textbf{The surrogate configurations compared during development}, as printed in
    \texttt{research/sim-data/quality-validation.md} \S11\docnum{} (an earlier snapshot of the data; MAE in
    $10^{-3}$ bpb). The first LOPO and the LOKVO columns use the pairs named in the simulator's specification;
    ``recipe'' is the subset that changes the recipe, and ``Decisive'' its decisive pairs as defined in that report.
    The default configuration, in bold, is the one used throughout the paper.}
  \label{tab:surrogate-configs}
  % GENERATED by paper/analysis -- surrogate configurations, research/sim-data/quality-validation.md section 11 (values as printed there); MAE in 1e-3 bpb (simulator.py)
{\footnotesize\setlength{\tabcolsep}{4pt}
\begin{tabular}{@{}lp{4.2cm}rrrrrrrr@{}}
\toprule
& & & \multicolumn{2}{c}{LOPO, \val{sc-spec-n} pairs} & \multicolumn{2}{c}{LOPO, recipe} & Decisive & \multicolumn{2}{c}{LOKVO, \val{sc-spec-n} pairs} \\
\cmidrule(lr){4-5}\cmidrule(lr){6-7}\cmidrule(lr){9-10}
Config & Change & Runs & Sign & MAE & Sign & MAE & sign & Sign & MAE \\
\midrule
{A0} & original conventions: K59 defaults for missing knobs, no leakage guard, every code version its own era & 131 & 84\% & 0.66 & 81\% & 0.68 & 92\% & 71\% & 0.86 \\
{A1} & code defaults for missing knobs, code lineage, leakage guard & 131 & 85\% & 0.38 & 82\% & 0.34 & 92\% & 63\% & 0.87 \\
{A2} & A1 + one era for M9 and its flag-gated successors & 131 & 81\% & 0.38 & 77\% & 0.34 & 90\% & 70\% & 0.73 \\
\textbf{A3} & A2 + bowl prior on tuned knobs, bowls $\ge 0$ (the default) & 131 & 85\% & 0.38 & 82\% & 0.34 & 95\% & 84\% & 0.60 \\
{A4} & A3, slope prior SD 0.002 & 131 & 86\% & 0.38 & 84\% & 0.34 & 97\% & 84\% & 0.61 \\
{A5} & A3, seed prior SD 0.0009 & 131 & 85\% & 0.38 & 82\% & 0.34 & 95\% & 84\% & 0.69 \\
{A6} & A3, era grouping off & 131 & 90\% & 0.38 & 89\% & 0.34 & 97\% & 78\% & 0.73 \\
{A7} & A3, K59 defaults for missing knobs & 131 & 74\% & 0.66 & 69\% & 0.68 & 87\% & 70\% & 0.81 \\
{A8} & A3, bowl prior mean 0.001 & 131 & 86\% & 0.42 & 84\% & 0.39 & 97\% & 85\% & 0.69 \\
{A9} & A3, knob prior SD 0.002 & 131 & 85\% & 0.37 & 82\% & 0.33 & 95\% & 84\% & 0.64 \\
{A10} & A3 without the bowl constraint & 131 & 85\% & 0.38 & 82\% & 0.34 & 95\% & 84\% & 0.60 \\
{A11} & A3 fitted on M9+ runs only & 71 & 85\% & 0.65 & 82\% & 0.65 & 97\% & 82\% & 0.86 \\
\bottomrule
\end{tabular}
}

\end{table}

\section{Seed sweeps}

\begin{table}[!htbp]
  \centering
  \caption{\textbf{Warm-up paths in three seed sweeps of a fixed recipe} (full chip runs that differ only in the
    seed; $10^{-3}$ bpb). ``Spike'': step-10 loss $\ge$\val{spike-l10} nats; ``Interm.'': between \val{l10-interm-lo} and
    \val{spike-l10}. The difference is a Welch mean difference with 95\% CI. Generated by
    \texttt{paper/analysis/seeds.py}.}
  \label{tab:seeds}
  % GENERATED by paper/analysis -- warm-up path per seed sweep, 1e-3 bpb (seeds.py)
{\footnotesize\setlength{\tabcolsep}{4pt}
\begin{tabular}{@{}p{4.0cm}rrrrlrr@{}}
\toprule
Seed sweep & Seeds & Curve & Spike & Interm. & Spike $-$ good [95\% CI] & SD all & SD good \\
\midrule
K57 recipe (M9), chip C & 13 & 13 & 4 & 1 & \ensuremath{+}0.93 [\ensuremath{+}0.31, \ensuremath{+}1.56] & 0.61 & 0.42 \\
LeakyReLU(0.35)$^2$ (M12), chip C & 9 & 7 & 2 & 0 & \ensuremath{+}1.52 [\ensuremath{-}0.06, \ensuremath{+}3.10] & 1.32 & 0.53 \\
LeakyReLU(0.35)$^2$ (M12), chip D & 7 & 7 & 4 & 0 & \ensuremath{+}1.38 [0.00, \ensuremath{+}2.77] & 0.95 & 0.66 \\
\midrule
pooled (centred on each group's good-path mean) & 29 & 27 & 10 & & \ensuremath{+}1.23 [\ensuremath{+}0.80, \ensuremath{+}1.66] & & \\
\bottomrule
\end{tabular}
}

\end{table}

\clearpage
\section{The contribution pipeline}
\label{app:contrib}

\paragraph{The record.} One JSON object per run with a published schema: hardware, time budget, an optional
published base recipe (\upl{K82s4}, \upl{K60}, \upl{K59}) plus only the knobs that were changed, seed, steps, the
local rehearsal score and/or the official score, and an explicit consent flag. A run can be submitted through the
local app, a GitHub issue form or the command line; the app and the command line only \emph{build} a pre-filled issue
link and never send anything themselves.

\paragraph{Validation.} A schema check run by a small built-in interpreter; a scan for anything that looks like a
secret or personal data (access keys, account and resource identifiers, tokens, private keys, e-mail and host IP
addresses), with the field named but the value never repeated; plausible ranges for every knob the model reads,
rejection of non-finite numbers, and a finite feature vector and prediction; duplicate detection by a content hash;
and an outlier check against the current model, which flags a record with $|z|>\val{guard-outlier-z}$ for the
reviewer.

\paragraph{Entering the fit.} Each (hardware, base recipe) combination gets its own fixed effect with the same prior
as a chip. Runs shorter than \val{guard-min-steps} steps are stored but never fitted. A published base whose newest
levers lie outside the data (\upl{K82s4}) is anchored to its measured rehearsal (\val{guard-anchor}) until its effect
has data of its own.

\paragraph{Review and merge.} Only a label added by an account with write access approves a record, and editing the
issue afterwards removes the label. The approval opens a pull request that appends the record exactly as approved, on a
branch named by its content hash. Merging triggers a guarded refit whose output arrives as a second pull request;
nothing pushes to the main branch. Workflows run with read-only tokens wherever possible, pin every action to a commit
hash and never execute contributed content.

\paragraph{Stress test.} \Cref{tab:guard} pushes synthetic records through the real validation and guard code in
memory (\texttt{paper/analysis/contrib\_guard.py}). Honest records are the model's own prediction plus a hardware
bias and seed-sized noise, on knob values the data have seen; they moved the base predictions by at most
\val{guard-honest-c-shift} and were accepted. A single record 0.01~bpb worse than predicted was absorbed by its own
hardware effect (base shift \val{guard-outlier-shift}) and was not flagged, because an unseen hardware effect widens
the predictive SD. The coordinated attack was refused at every size tried, down to two records, which would have moved a
published base prediction by \val{guard-attack-two-shift}.

\begin{table}[!htbp]
  \centering
  \caption{\textbf{Refit-guard stress test.} Changes in $10^{-3}$ bpb: built-in in-sample MAE, pair-validation MAE
    (leave one pair out), and the largest move of a published base recipe's prediction. ``Flagged'': records with
    $|z|>\val{guard-outlier-z}$. The guard refuses a refit when either MAE worsens by more than
    \val{guard-max-worse} or a base prediction moves by more than \val{guard-max-shift}.}
  \label{tab:guard}
  % GENERATED by paper/analysis -- refit-guard stress test, all deltas x 1e-3 bpb (contrib_guard.py)
{\footnotesize
\begin{tabular}{@{}p{4.6cm}rrrrrl@{}}
\toprule
Scenario & Records & Flagged & $\Delta$ in-sample & $\Delta$ pairs & Base shift & Guard \\
\midrule
honest, K60 base, 8 runs, bias +0.0015 & 8 & 0 & \ensuremath{+}0.002 & \ensuremath{-}0.025 & 0.01 & accepted \\
honest, K82s4 base, 8 runs, bias $-$0.0010 & 8 & 0 & \ensuremath{+}0.001 & \ensuremath{-}0.025 & 0.08 & accepted \\
honest, both bases, 24 runs, bias +0.0015 & 24 & 0 & \ensuremath{+}0.001 & \ensuremath{-}0.030 & 0.23 & accepted \\
one outlier (+0.01 bpb) & 1 & 0 & 0.000 & \ensuremath{-}0.001 & 0.09 & accepted \\
slope attack, 2 sock puppets & 2 & 1 & \ensuremath{+}0.128 & \ensuremath{+}0.022 & 4.73 & \textbf{refused} \\
slope attack, 4 sock puppets & 4 & 2 & \ensuremath{+}0.133 & \ensuremath{+}0.027 & 5.11 & \textbf{refused} \\
slope attack, 8 sock puppets & 8 & 4 & \ensuremath{+}0.139 & \ensuremath{+}0.028 & 5.42 & \textbf{refused} \\
\bottomrule
\end{tabular}
}

\end{table}

\paragraph{How much the stress test shows.} The stress test is weak evidence and we present it as such. Its
``honest'' records are generated from the model's own predictions plus noise, so accepting them is close to
automatic; its attack, a coordinated pair of sock-puppet records that pull a shared coefficient, was designed by us
with the thresholds known, and there is no adaptive attacker. The guard refused that attack down to two records and
accepted the honest ones. An attacker who knows the thresholds can stay under them, and per-hardware effects can
absorb slow drifts. What the guard does is turn the cheapest damaging attack, pulling a shared coefficient, into a
visible failure.

\clearpage
\section{The final recipe}
\label{app:recipe}

\begin{table}[!htbp]
  \centering
  \caption{\textbf{The final recipe, \upl{K82s4}.} About \val{params}\docnum{} parameters; every value is a default
    of \texttt{recipes/K82s4/train.py} (knob settings in \cref{tab:hyper}).}
  \label{tab:recipe}
  \small
  \begin{tabular}{@{}p{2.2cm}p{13.6cm}@{}}
    \toprule
    Component & Setting \\
    \midrule
    Model & 9 blocks, width 1{,}024; four 256-wide query heads sharing one key/value head (GQA); value embeddings
      fed to blocks 0 and 8; LeakyReLU(0.35)$^2$ MLP ($4\times$); rotary embeddings with a key offset (radius 0.3);
      blocks 6--8 reuse block~5's attention source (\knob{FF_ATTN_SRC=5:6,7,8}); logit soft cap 15. \\
    Objective & Next-token loss plus a 3-token multi-token-prediction loss with weights 1/0.5/0.25, faded to plain
      next-token prediction in eight stages by 45\% of the run. \\
    Optimizers & Muon (5 Newton--Schulz iterations, fused update) for matrices; AdamW for embeddings, scalars and
      the head; MLP output projection at learning rate $\times0.8$. \\
    Schedule & Charged-time target 1{,}795~s; batch warm-up 65k $\to$ 131k $\to$ 262k tokens per step
      (\knob{FF_ACCUM_SCHED=1:0.06,2:0.2,4}) with learning rate $\times0.35$ and $\times0.6$ in the small-batch phases;
      20 warm-up steps; linear cooldown over the last 70\% of the run. \\
    Data & Stream-row loader; from micro-batch \val{rowpool-from} on, rows are drawn without replacement from a pool
      of \val{rowpool-size} micro-batches (shuffle salt 4). \\
    Final weights & 0.6 blend of live weights and an EMA updated every 32 steps; the EMA update is compiled
      during the excluded start-up (\knob{FF_EMA_PREWARM}). \\
    Seed & 73. \\
    \bottomrule
  \end{tabular}
\end{table}

\Cref{tab:hyper} lists the settings of \upl{K82s4} that the paper refers to, read from
\path{ffsim/examples/recipe-K82s4.json}, which records all \val{k82-nknobs} knobs of
\path{recipes/K82s4/train.py}. Most of the rest are switches for mechanisms that are off, or for profiling. The
shuffle salt is a code constant, not a knob: \upl{K82s7} differs from \upl{K82s4} in that one line, and so is a
different file with a different hash.

\begin{table}[!htbp]
  \centering
  \caption{\textbf{Settings of the final recipe \upl{K82s4}.}}
  \label{tab:hyper}
  % GENERATED by paper/analysis -- ffsim/examples/recipe-K82s4.json (make_all.py)
{\footnotesize
\begin{tabular}{@{}lp{3.2cm}p{7.6cm}@{}}
\toprule
Knob & Value & Meaning \\
\midrule
\texttt{FF\_DEPTH} & \texttt{9} & transformer blocks \\
\texttt{FF\_ASPECT\_RATIO} & \texttt{113} & width = aspect $\times$ depth, rounded to the head size (1{,}024) \\
\texttt{FF\_HEAD\_DIM} & \texttt{256} & query head size (four query heads) \\
\texttt{FF\_KV\_HEADS} & \texttt{1} & key/value heads (grouped-query attention) \\
\texttt{FF\_MLP\_MULT} & \texttt{4} & MLP expansion \\
\texttt{FF\_LEAKY\_RELU2} & \texttt{0.35} & LeakyReLU$^2$ negative slope \\
\texttt{FF\_VALUE\_EMBED} & \texttt{1} & value embeddings (fed to blocks 0 and 8) \\
\texttt{FF\_ATTN\_SRC} & \texttt{5:6,7,8} & attention-source reuse: blocks 6, 7, 8 reuse block 5 \\
\texttt{FF\_KEY\_OFFSET\_RAD} & \texttt{0.3} & rotary key-offset radius \\
\texttt{FF\_ROPE\_BASE} & \texttt{100000} & RoPE base \\
\texttt{FF\_SOFTCAP} & \texttt{15} & logit soft cap \\
\texttt{FF\_TRAIN\_SEQ} & \texttt{1024} & sequence length (tokens) \\
\texttt{FF\_MTP\_W} & \texttt{1,0.5,0.25} & multi-token prediction weights (next, +2, +3) \\
\texttt{FF\_MTP\_PHASES} & \texttt{0.2,0.45} & MTP fade: start and end as fractions of the run \\
\texttt{FF\_TOTAL\_BATCH} & \texttt{262144} & final batch (tokens per optimizer step) \\
\texttt{FF\_MB} & \texttt{8} & micro-batch (sequences per rank) \\
\texttt{FF\_ACCUM\_SCHED} & \texttt{1:0.06,2:0.2,4} & batch warm-up: $k$ micro-batches until a time fraction \\
\texttt{FF\_ACCUM\_LR} & \texttt{0.35,0.6} & LR multipliers of the two small-batch phases \\
\texttt{FF\_TIME\_TARGET} & \texttt{1795} & charged-time target (s) \\
\texttt{FF\_WARMUP} & \texttt{20} & LR warm-up steps \\
\texttt{FF\_COOLDOWN\_FRAC} & \texttt{0.7} & linear cooldown, fraction of the run \\
\texttt{FF\_MATRIX\_LR\_SCALE} & \texttt{2} & Muon LR scale (matrices) \\
\texttt{FF\_ADAMW\_LR\_SCALE} & \texttt{1.4142} & AdamW LR scale (embeddings, scalars, head) \\
\texttt{FF\_EMB\_LR} & \texttt{0.3} & embedding LR \\
\texttt{FF\_UNEMB\_LR} & \texttt{0.006} & output-head LR \\
\texttt{FF\_SCALAR\_LR} & \texttt{0.2} & scalar-parameter LR \\
\texttt{FF\_MLP\_PROJ\_LR} & \texttt{0.8} & MLP output-projection LR multiplier \\
\texttt{FF\_WD} & \texttt{0.02} & weight decay \\
\texttt{FF\_MOM\_PEAK} & \texttt{0.95} & Muon momentum peak \\
\texttt{FF\_NS\_STEPS} & \texttt{5} & Newton--Schulz iterations \\
\texttt{FF\_OPT\_FUSE} & \texttt{3} & fused optimizer update level \\
\texttt{FF\_XU\_QUEUE} & \texttt{48} & execution-queue depth \\
\texttt{FF\_ROW\_SHUFFLE} & \texttt{256} & row pool size (micro-batches) \\
\texttt{FF\_ROW\_SHUFFLE\_FROM} & \texttt{96} & row pool starts at this micro-batch \\
\texttt{FF\_EMA\_EVERY} & \texttt{32} & EMA update interval (steps) \\
\texttt{FF\_EMA\_BLEND} & \texttt{0.6} & final weights: blend of live and EMA weights \\
\texttt{FF\_EMA\_PREWARM} & \texttt{1} & compile the EMA update during excluded startup \\
\texttt{FF\_SEED} & \texttt{73} & training seed \\
\bottomrule
\end{tabular}
}

\end{table}

\section{Derivations}
\label{app:derivations}

\paragraph{Step equivalents.} At the price $\kappa$ of \cref{eq:exchange}, a change $\Delta$ is worth the same as
multiplying the step count by $\exp(|\Delta|/\kappa)$; the step equivalent in \cref{tab:levers} is that factor minus
one.

\paragraph{Compute needed for a gap.} Inverting \cref{eq:exchange}, closing a gap $g$ needs a compute ratio of
$\exp(g/\kappa)$ (\cref{tab:gap}).

\paragraph{Pair noise floor.} If each run's score has independent residual noise of SD $\sigma_q$, the difference of
two single runs has SD $\sqrt2\,\sigma_q$ and mean absolute value $\sqrt2\,\sigma_q\sqrt{2/\pi}$, which is
\val{pair-noise-floor} at $\hat\sigma_q=\val{q-sigma}$ and \val{pair-noise-floor-nofloor} at the unfloored estimate. No
pair predictor can have a lower expected MAE on single-run pairs.

\paragraph{Resolution of a difference.} If two uploads' oracle errors are independent with SD $s$, their difference
has SD $\sqrt2\,s$ and a 95\% half-width of $1.96\sqrt2\,s$ (\cref{sec:resolution}).

\paragraph{The winner's curse for salts.} Write a salt's rehearsal advantage as $a = \theta + \eta$, with true effect
$\theta \sim \mathcal{N}(0, s_\theta^2)$ and a part $\eta\sim\mathcal{N}(0, s_\eta^2)$ that does not transfer to the
official run. Then $\mathbb{E}[\theta \mid a] = \lambda a$ with $\lambda = s_\theta^2/(s_\theta^2 + s_\eta^2)$. What
is measured is the spread of rehearsal scores, $s_a^2 = s_\theta^2 + s_\eta^2$, so $\lambda = 1 - s_\eta^2/s_a^2$,
with $s_a$ the salt SD (\val{salt-sd}\docnum{}). Noise that enters only the official run is independent of $a$ and
does not enter $s_\eta$. On one chip, $s_\eta^2$ is the variance of the text gap (SD \val{gap-sd}) plus that of the
rehearsal's step-count jitter, $\kappa$ times the within-sweep SD of steps over their mean
(\val{wc-jit-lo}--\val{wc-jit-hi}), so $s_\eta=\val{wc-eta-lo}$--\val{wc-eta-hi} and
$\lambda=\val{wc-lam-lo}$--\val{wc-lam-hi} over the ranges of $s_a$ and the jitter. When the draws ran on different
chips, a between-chip offset SD $s_b=\val{wc-between}$ (method of moments over the uploads since \upl{K50}) adds to
both the spread and the part that does not transfer, so the share lost is
$(s_\eta^2+s_b^2)/(s_a^2+s_b^2)=\val{wc-shrink-cross-lo}$--\val{wc-shrink-cross-hi} (\cref{sec:salts}). This
$s_\eta$ assumes that the rehearsal itself is repeatable apart from its step jitter: the remainder of the within-chip
offset SD (\val{offdec-rest}, \cref{app:offset-detail}) is treated as official-side. If part of it is rehearsal
noise, $s_\eta$ is larger and $\lambda$ smaller; repeated cold runs of one file and seed would measure that
noise directly.

\paragraph{Prediction interval for the offset.} Under \cref{eq:offset} with $\gamma\equiv0$, and $\hat\mu$ and $s$
estimated from $n$ uploads, $(y_{\text{new}} - r_{\text{new}} - \hat\mu)/(s\sqrt{1+1/n})$ follows a Student $t$
distribution with $n-1$ degrees of freedom, which gives the bands of \cref{fig:offset-prospective} and the coverage
levels of \cref{tab:coverage}.

\section{Numbers quoted from campaign documents}
\label{app:docfacts}

Numbers marked \docnum{} in the text are quoted from the campaign documents in \texttt{docs/} (and two validation
reports in \texttt{research/sim-data/}) because the records behind them are not released; a few non-numeric facts
marked the same way come from the campaign's unreleased log (``campaign log'' in the table). They are registered once in
\texttt{paper/analysis/docfacts.py}, \texttt{exchange.py} and \texttt{simulator.py}. \Cref{tab:docfacts} lists them;
``contest rules'' there refers to the published rules \citep{trainiumfrontier2026rules}.

{\footnotesize
% GENERATED by paper/analysis/make_all.py -- values quoted from campaign documents
\begin{longtable}{@{}p{3.6cm}p{12.0cm}@{}}
\caption{\textbf{Values quoted from campaign documents} (marked \docnum{} in the text).}\label{tab:docfacts}\\
\toprule
Value & Source (document, section: what) \\
\midrule
\endfirsthead
\toprule
Value & Source (continued) \\
\midrule
\endhead
\bottomrule
\endlastfoot
1{,}800 & contest rules (trainiumfrontier2026rules): 30 minutes excluding start-up/compilation \\
6--12~December 2026 & contest rules: NeurIPS 2026 competition event \\
7~October to 4~November 2026 & contest rules \\
11:59~PM PT on 30~September 2026 & contest rules: Round 1 deadline \\
about 20M & contest rules: pinned validation shard, \textasciitilde{}20M tokens \\
7--12\% & docs/EXACT-ORACLE.md traps \\
3.5\% & docs/EXACT-ORACLE.md traps \\
1.041525 & docs/EXACT-ORACLE.md why it works (19 Sep) \\
5\ensuremath{\times}10\textsuperscript{\ensuremath{-}6} & docs/EXACT-ORACLE.md \\
1.041520 & docs/EXACT-ORACLE.md why it works \\
0.0002--0.0006 & docs/EXACT-ORACLE.md traps \\
2{,}097{,}152 & docs/EXACT-ORACLE.md: first public tokens used by the rehearsal \\
6--8 & docs/EXACT-ORACLE.md traps: seconds of headroom under the 1,800 s cap \\
0.9619 & docs/EXACT-ORACLE.md: largest miss against a written projection \\
\ensuremath{+}0.0069 to \ensuremath{+}0.0070 & docs/EXACT-ORACLE.md reproducing it: final-day offset \\
0.9614--0.9615 & docs/EXACT-ORACLE.md reproducing it \\
0.9611--0.9619 & docs/EXACT-ORACLE.md reproducing it \\
\ensuremath{+}0.0014 & docs/EXACT-ORACLE.md: organiser-side K44 rerun \\
2.3\% & docs/EXACT-ORACLE.md: organiser-side K44 rerun \\
\ensuremath{+}0.0070 & docs/EXACT-ORACLE.md: 20M vs first 2M public tokens \\
\ensuremath{-}0.00149 & docs/FINDINGS.md s3, mean of four pairs \\
\ensuremath{-}0.00181, \ensuremath{-}0.00136, \ensuremath{-}0.00106, \ensuremath{-}0.00173 & docs/FINDINGS.md s3 \\
0.1\% & docs/FINDINGS.md s3 \\
\ensuremath{-}0.0006 & docs/FINDINGS.md s3, three chip pairs \\
\ensuremath{-}0.0004 & docs/FINDINGS.md s3 \\
6--11\% & docs/FINDINGS.md s5 \\
10.6\% & docs/FINDINGS.md s5 \\
\ensuremath{-}0.0005 & docs/FINDINGS.md s3, one chip pair \\
15--30 & docs/FINDINGS.md s3 / docs/EXACT-ORACLE.md traps \\
430 & docs/FINDINGS.md s4 \\
0.9614 & docs/FINDINGS.md s1: K82s7 projected about 0.9614 \\
1:41~AM CDT on 1~October, i.e.\ 11:41~PM PT on 30~September & docs/FINDINGS.md s1 \\
\ensuremath{-}0.0017 & docs/FINDINGS.md s1 \\
0.0013 & docs/FINDINGS.md s1: each structural change worth 0.0013 or more \\
6 & docs/FINDINGS.md s4: 1 of about 6 new seeds landed on the good path \\
0.8\% & docs/FINDINGS.md s3 \\
124M & docs/FINDINGS.md s2 \\
1{,}024 & docs/FINDINGS.md s3: a pool of 1,024 gave the same result as 256 \\
64 & docs/FINDINGS.md s3: a pool of 64 helped less \\
\ensuremath{-}0.0023 & docs/FINDINGS.md s3: chip E pair \\
0.959523 & docs/FINDINGS.md s3 \\
0.957259 & docs/FINDINGS.md s3 \\
96 & docs/FINDINGS.md s2 (FF\_ROW\_SHUFFLE\_FROM) \\
0.0003 & docs/FINDINGS.md s5: other order tricks 0 +- 0.0003 \\
256 & docs/FINDINGS.md s2 (FF\_ROW\_SHUFFLE) \\
0.5\% & docs/FINDINGS.md s3 \\
\ensuremath{-}0.0026 & docs/FINDINGS.md s3 \\
2{,}329 & docs/FINDINGS.md s3 \\
2{,}318 & docs/FINDINGS.md s3 \\
0.0003--0.0004 & docs/FINDINGS.md s3 \\
0.0008 & docs/FINDINGS.md s3: four salts spanned 0.0008 on chip \\
61 & docs/FINDINGS.md s5 \\
0.0010--0.0012 & docs/FINDINGS.md s4 \\
40\% & docs/FINDINGS.md s4 \\
\ensuremath{+}0.003 to \ensuremath{+}0.0045 & docs/FINDINGS.md s4 \\
2.8\% & docs/FINDINGS.md s6 \\
2.6\% & docs/FINDINGS.md s6 \\
0.7\% & docs/FINDINGS.md s6 \\
9.1\% & docs/FINDINGS.md s6 \\
24~September & docs/GAP-ANALYSIS.md \\
19\% & docs/GAP-ANALYSIS.md: idle on host dispatch \\
0.0005 & docs/GAP-ANALYSIS.md: recipe changes on the final day \\
7:40~AM CDT on 1~October & docs/GAP-ANALYSIS.md: final snapshot \\
30\% & docs/GAP-ANALYSIS.md: roughly 30\% MFU \\
285k & docs/GAP-ANALYSIS.md where the compute goes \\
0.00056 & docs/SIMULATOR.md \\
\ensuremath{+}0.09 and \ensuremath{-}0.37$\times10^{-3}$ & docs/SIMULATOR.md route 2: cooldown 0.7 chip pairs \\
2.5 & docs/SIMULATOR.md \\
185 & docs/SIMULATOR.md \\
1.35 & docs/SIMULATOR.md \\
32 & docs/SIMULATOR.md \\
104k & docs/SIMULATOR.md \\
76k & docs/SIMULATOR.md \\
\ensuremath{-}0.00105 & docs/SIMULATOR.md route 2: chip C pair vs ReLU\textsuperscript{2} \\
282 & docs/SIMULATOR.md: local candidates around K60 \\
282 & docs/SIMULATOR.md \\
7.6\% & docs/SIMULATOR.md route 1: a 96-parameter gate \\
3\% & docs/SIMULATOR.md route 1: unknown-mechanism prior \\
1.8\% & docs/SIMULATOR.md route 1: one fewer elementwise pass \\
3.2\% & docs/SIMULATOR.md route 1: QKV fusion slower \\
6--10 & docs/SIMULATOR.md \\
0.3\% & docs/SIMULATOR.md \\
0.0005 & docs/simulator-report-20260929.md s5: attention-source pairs over-predicted by about 0.0005 \\
\ensuremath{+}0.00115 & docs/simulator-report-20260929.md section 5: AS-a s97 (D), raw \\
41 & docs/simulator-report-20260929.md section 5: 2109 vs 2150 steps \\
15 & docs/simulator-report-20260929.md section 5, scored in simulator.py \\
0.095--0.115 & ffsim/quality.py docstring: slope magnitude at about 900 steps \\
0.00005 & official scores are published to 4 decimals (results CSV) \\
1{,}188 & research/sim-data/quality-validation.md header: run records when the report was written \\
136 & research/sim-data/simulate-validation.md section 2 \\
5.5 & research/sim-data/simulate-validation.md section 2 \\
\end{longtable}

}

\section{Notation}
\label{app:notation}

\begin{table}[!htbp]
  \centering
  \caption{\textbf{Notation.}}
  \label{tab:notation}
  \small
  \begin{tabular}{@{}lp{12.4cm}@{}}
    \toprule
    Symbol & Meaning \\
    \midrule
    $S$ & optimizer steps in the budget \\
    $k$ & micro-batches (65{,}536 tokens each) accumulated per optimizer step; $\bar t_k$ the step time of that phase \\
    $\kappa$, $\rho$ & the price of a step: $\kappa=\val{kappa}$ bpb per unit $\ln S$; $\rho=\kappa/100$ per 1\% of steps (\cref{eq:exchange}) \\
    $r_i$, $y_i$ & rehearsal and official score of upload $i$ \\
    $\mu$, $\gamma_c$ & text part and chip part of the offset (\cref{eq:offset}) \\
    $\sigma_{\text{off}}$, $s$ & SD of the offset residual, and its sample estimate \\
    $T$, $t_{\text{res}}$ & charged-time target and the stop reserve (\cref{eq:decomp}) \\
    $q$, $L$ & rehearsal score predicted by the surrogate, and $L=\ln(S/2300)$ (\cref{eq:quality}) \\
    $\alpha$, $\delta$, $b$, $c$, $\beta_j$, $u$ & surrogate era, chip, step, curvature, knob and seed effects
      (\cref{eq:quality}) \\
    $\sigma_q$, $s_j$, $m$ & surrogate residual SD, prior SDs and prior means \\
    $\varepsilon_i$ & residual of the offset model for upload $i$ (\cref{eq:offset}) \\
    $\varepsilon_q$ & residual of the quality surrogate, SD $\sigma_q$ (\cref{eq:quality}) \\
    $\lambda$, $s_a$, $s_\theta$, $s_\eta$, $s_b$ & winner's-curse shrinkage, the spread of rehearsal advantages, its
      transferable and non-transferable parts, and the between-chip offset SD (\cref{app:derivations}) \\
    \bottomrule
  \end{tabular}
\end{table}
